\chapter*{前言}
\markboth{前言}{前言}
\addcontentsline{toc}{chapter}{前言}
\SourcePage{5}
Fukaya 范畴素有难以入门之名。这既是因为学习它需要大量背景知识——来自同调代数、辛几何和几何分析——也是因为它的基本定义本身颇为复杂。本书希望为有意了解 Fukaya 范畴、并可能将其用于自身工作的研究生和研究人员提供一份参考资料。我尽量把注意力集中在一组相当基础的课题上，力求尽可能精确地加以叙述，并补上某些早期文献中的空缺。这使本书的风格颇为简峭（因此，深入研读本书时，最好同时阅读一些讨论应用的论文）。本书的第二个目标，是介绍若干此前尚未发表的结果；这些结果把 Fukaya 范畴与 Lefschetz 纤维化理论联系起来。这一目标在最后几节中逐渐占据主导地位，正文也随之逐步转变为一部研究专著。

我广泛借鉴了许多人的工作，其中首先要提到的是 Fukaya、Kontsevich 和 Donaldson。Fukaya 的奠基性贡献，当然是把 $\Ai$ 结构引入辛几何。在代数方面，他率先使用了 Yoneda 嵌入的 $\Ai$ 版本，而我们将在全书中系统地采用这一工具。此外，若干几何思想，例如 Pin 结构的作用，以及借助带参数模空间构造 $\Ai$ 同态的方法，都来自 Fukaya、Oh、Ohta 和 Ono 的工作。Kontsevich 引入了 $\Ai$ 范畴的导出范畴，并建立了其中的大部分理论，特别是正合三角的内在刻画。他还猜想了 Dehn 扭转与扭转函子之间的关系，而这正是本书的主要结果之一。最后，在与 Barannikov 的合作中，他提出了 Lefschetz 纤维化的 Fukaya 范畴的一种构造；我们采用的定义表面上有所不同，但想来应当与之等价。Donaldson 的影响同样贯穿始终。除了在 Lefschetz 束方面的开创性工作之外，他还引入了匹配圈，并提出以此为出发点建立 Floer 上同调的组合公式；这一设想在本书中确实得到了部分实现。其他数学家也作出了重要贡献。例如，我们对 Picard–Lefschetz 理论的部分叙述反映了 Auroux 的观点。Smith 的一个结果——四维 Lefschetz 束中的消失圈必然充满纤维——对于提出这些圈可能“分裂生成”Fukaya 范畴的设想起了关键作用。此外，Fukaya–Smith 关于余切丛的工作，为有关 Lefschetz 纤维化的一些更大胆的想法提供了很好的试验场。我们处理横截性问题的方法，源于与 Lazzarini 的数次交谈。最后，Abouzaid 的建议大大改进了关于辛嵌入的讨论。

\SourcePage{6}
在机构方面，我要感谢 ETH Zürich 邀请我讲授作为本书基础的课程；感谢 NSF 提供经费支持（项目编号 DMS-9729992 和 DMS-0405516）；还要感谢 EMS Publishing House 接受书稿，并悉心推动本书出版。
